\section{The derivative and its primitive recursive evaluator}
\label{sec:derivative}

\subsection{A finite construction of the inverse matrix}

\begin{lemma}[Neumann inverse with a rate]\label{lem:neumann}
If \(V\) is a represented \(m\times m\) matrix and
\(\norm{\Id-V}\le1/2\), write \(E=\Id-V\) and
\[
 S_N(E)=\sum_{j=0}^{N-1}E^j,\qquad S_0=0.
\]
There is a represented inverse \(S=V^{-1}\), with
\[
 \norm S\le2,\qquad
 \norm{S-S_N(E)}\le2^{1-N}.
 \tag{Neumann tail}
\]
Both inverse identities hold.
\end{lemma}

\begin{proof}
Finite multiplication and telescoping give
\begin{equation}\label{eq:geomidentity}
 (\Id-E)S_N(E)=S_N(E)(\Id-E)=\Id-E^N.
\end{equation}
The finite norm estimates are
\[
 \norm{E^N}\le2^{-N},\qquad
 \norm{S_{N+L}(E)-S_N(E)}
 \le\sum_{j=N}^{N+L-1}2^{-j}\le2^{1-N}.
\]
They produce a regular name for \(S\) after an explicit index
shift and rational evaluation of the finitely many matrix entries.
The inverse identities follow by observing \eqref{eq:geomidentity}
at any requested precision with sufficiently large \(N\).
The norm bound follows from the finite sum bound \(2\).

For the object-theory proof, do not induct on the real-matrix
inequality \(\norm{E^N}\le2^{-N}\) as an infinite-precision
statement. Fix the target \(N\) and observation precision first;
use a rational approximation to \(E\), the finite polynomial
identity, and the perturbation estimate below. The same estimates
also give an entirely rational evaluator for \(S\). The matrix
loop identities themselves are quantified finite arithmetic
identities proved by quantifier-free induction on coded arrays.
\end{proof}

\subsection{Deriving differentiability from the Taylor contract}

At a fixed parameter \(x\), abbreviate
\[
 U_x=U(x,g(x)),\quad V_x=V(x,g(x)),\quad
 D_x=-V_x^{-1}U_x.
\]
For \(x,x+h\in X\), put
\(\Delta=g(x+h)-g(x)\). The Lipschitz bound gives
\[
 \norm{(h,\Delta)}
 \le C\norm h,\qquad C=\max(1,2M).
\]
Let \(c=\lgup(C)\), and suppose
\(\norm h\le2^{-\mu(p+c+1)-c}\).
Then \(\norm{(h,\Delta)}\le2^{-\mu(p+c+1)}\).
Apply \eqref{eq:taylor} to the two points on the graph:
\[
 0=H(x+h,g(x+h))-H(x,g(x))
   =U_xh+V_x\Delta+R,
\]
where
\[
 \norm R\le2^{-(p+c+1)}C\norm h.
\]
Since \(\norm{V_x^{-1}}\le2\),
\begin{align*}
 \norm{\Delta-D_xh}
 &=\norm{-V_x^{-1}R}\\
 &\le 2^{-(p+c)}C\norm h
 \le2^{-p}\norm h.
\end{align*}
This proves both differentiability and the modulus \(\mu_g\) in
\eqref{eq:outputmoduli}. No division by \(\norm h\) is used;
the argument includes \(h=0\).
The identity defining \(D_x\) also yields \(\norm{D_x}\le2M\).

\subsection{Continuity of the derivative}

At \(x'\), write \(U_{x'},V_{x'}\) analogously.
Finite algebra gives the inverse difference identity
\[
 V_{x'}^{-1}-V_x^{-1}
   =V_{x'}^{-1}(V_x-V_{x'})V_x^{-1}.
\]
If each of the two derivative blocks changes by at most \(\delta\),
then
\[
 \norm{V_{x'}^{-1}-V_x^{-1}}\le4\delta
\]
and
\begin{equation}\label{eq:differenceD}
 \norm{D_{x'}-D_x}
 \le 2\delta+4M\delta=(2+4M)\delta.
\end{equation}
The distance between the graph points is at most
\(C\norm{x'-x}\). Put \(d=\lgup(2+4M)\).
If \(\norm{x'-x}\le2^{-\nu(p+d)-c}\),
\eqref{eq:duc} supplies \(\delta=2^{-(p+d)}\).
Equation~\eqref{eq:differenceD} gives the modulus \(\nu_g\).
This completes the \(C^1\) part of Theorem~\ref{thm:main}.

\subsection{Rational computation of the derivative}

The formula involving \(V_x^{-1}\) is not an instruction to
search for an inverse of an unknown real matrix. Here is a
bounded rational algorithm.

\begin{lemma}[Coefficient errors in a finite Neumann sum]
\label{lem:perturb}
Suppose \(\norm E\le1/2\), \(\norm U\le M\), and rational
approximations satisfy
\[
 \norm{\widetilde E-E}\le\delta,\qquad
 \norm{\widetilde U-U}\le\delta,\qquad \delta\le1/4.
\]
Then, for every \(N\),
\begin{align}
 \norm{S_N(\widetilde E)}&\le4,\label{eq:snorm}\\
 \norm{S_N(\widetilde E)-S_N(E)}&\le8\delta,\label{eq:sperturb}\\
 \norm{S_N(\widetilde E)\widetilde U-S_N(E)U}
 &\le(4+8M)\delta.\label{eq:productperturb}
\end{align}
\end{lemma}

\begin{proof}
We have \(\norm{\widetilde E}\le3/4\); the finite geometric
sum proves \eqref{eq:snorm}. For powers, use the exact telescoping
identity
\[
 \widetilde E^j-E^j
 =\sum_{i=0}^{j-1}
   \widetilde E^i(\widetilde E-E)E^{j-1-i}.
\]
Summing these finite inequalities for \(j<N\), the sum over the
triangular index set is bounded by the product of two finite
geometric sums:
\[
 \norm{S_N(\widetilde E)-S_N(E)}
 \le \delta
  \left(\sum_{i=0}^{N-1}(3/4)^i\right)
  \left(\sum_{l=0}^{N-1}(1/2)^l\right)
 \le8\delta.
\]
Finally split the product difference into
\[
 S_N(\widetilde E)(\widetilde U-U)
 +(S_N(\widetilde E)-S_N(E))U
\]
and apply the first two bounds.
\end{proof}

Choose once and for all
\[
 c_D=\max(3,\lgup(8+20M)).
\]
For requested derivative precision \(p\), set
\[
 q=p+c_D,\qquad N=q,\qquad j=\nu(q+1).
\]
For rational \(x\in X\), compute \(y=\Ghat(x,j)\) and evaluate
\(\Jhat((x,y),q+1)\). Let its blocks be
\(\widetilde U,\widetilde V\), and put
\(\widetilde E=\Id-\widetilde V\). The bound
\eqref{eq:representation} is at most \(2^{-j}\); graph-point
continuity and evaluation each contribute at most \(2^{-(q+1)}\).
Consequently both matrix errors are at most
\(\delta=2^{-q}\). Return
\begin{equation}\label{eq:derivativeprogram}
 \widehat D(x,p)
   =-S_q(\widetilde E)\widetilde U.
\end{equation}
All arithmetic after obtaining the matrix coefficients is exact
rational arithmetic. Lemma~\ref{lem:perturb} and the true
Neumann tail give
\begin{align*}
 \norm{\widehat D(x,p)-Dg(x)}
 &\le(4+8M)2^{-q}+2M\,2^{-q}\\
 &=(4+10M)2^{-(p+c_D)}
 \le2^{-(p+1)}.
\end{align*}
The evaluator is therefore primitive recursive. At a named
parameter it extends using the already proved modulus \(\nu_g\).

For completeness, the analogous inverse-only evaluator uses a
rational approximation to \(E\) of accuracy \(\delta\) and returns
\(S_N(\widetilde E)\); its error is at most
\(8\delta+2^{1-N}\). For example \(N=p+3\) and
\(\delta=2^{-(p+5)}\) give an error at most \(2^{-(p+1)}\).
Thus even the auxiliary inverse name has a direct finite
production rule.

\subsection{Why the proof remains first order}

Symbols such as \(V_x^{-1}\) in the human proof abbreviate the
specific name just constructed. Matrix multiplication is a finite
fold, and the geometric and power-difference identities have
explicit finite loop proofs. To establish an inverse identity at
precision \(p\), instantiate these finite identities at the
displayed \(N,\delta\); there is no invocation of completeness of
an unspecified Banach space and no choice of a matrix inverse
from a nonempty set. The formalization must preserve this
finite-observation reading, rather than simply import a classical
real inverse theorem.
